\documentclass[12pt,a4paper]{article}

\usepackage[utf8]{inputenc}
\usepackage[T1]{fontenc}
\usepackage{geometry}
\usepackage{hyperref}
\usepackage{xcolor}
\usepackage{listings}

\geometry{margin=2.5cm}

\lstdefinestyle{code}{
    basicstyle=\ttfamily\small,
    breaklines=true,
    frame=single,
    columns=fullflexible
}

\title{README \\ T17 - Analiza Semantica Cod Sursa}
\author{Proiect PCD}
\date{2026}

\begin{document}

\maketitle
\tableofcontents
\newpage

\section{Descriere proiect}

Proiectul T17 - Analiza Semantica Cod Sursa este o aplicatie client-server pentru analiza statica a fisierelor C. Clientii pot incarca fisiere sursa catre server, iar serverul ruleaza un analyzer bazat pe libclang si genereaza rapoarte text.

Aplicatia include:

\begin{itemize}
    \item server TCP INET pentru clientii normali;
    \item socket UNIX pentru clientul admin;
    \item client C normal;
    \item client C admin;
    \item client Python;
    \item analyzer bazat pe libclang;
    \item autentificare simpla;
    \item upload si download de fisiere;
    \item rapoarte generate automat;
    \item loguri si statistici;
    \item job tracking;
    \item monitorizare \texttt{reports/} cu inotify.
\end{itemize}

\section{Tehnologii folosite}

\begin{itemize}
    \item C;
    \item Python;
    \item TCP sockets;
    \item UNIX domain sockets;
    \item \texttt{select()};
    \item \texttt{fork()};
    \item \texttt{exec()};
    \item \texttt{pipe()};
    \item \texttt{inotify};
    \item libclang;
    \item libconfig;
    \item Makefile.
\end{itemize}

\section{Structura proiectului}

\begin{lstlisting}[style=code]
Proiect_PCD/
|-- Makefile
|-- README.md
|-- README.tex
|-- config/
|   |-- app.cfg
|   |-- users.cfg
|-- docs/
|   |-- SRS_SDD.tex
|   |-- protocol.tex
|   |-- openapi.yaml
|-- python_client/
|   |-- client.py
|-- src/
|   |-- analyzer.c
|   |-- main.c
|   |-- server.c
|   |-- client.c
|   |-- admin_client.c
|   |-- watch_reports.c
|-- tests/
|   |-- sample.c
|   |-- bad.c
|   |-- bad_syntax.c
|   |-- bad_type.c
|   |-- bad_uninitialized.c
|   |-- good_math.c
|   |-- good_nested.c
|   |-- good_while.c
\end{lstlisting}

\section{Instalare dependinte}

Pe Ubuntu/WSL se pot instala dependintele cu:

\begin{lstlisting}[style=code]
sudo apt update
sudo apt install build-essential libclang-dev libconfig-dev python3
\end{lstlisting}

In functie de versiunea LLVM instalata, Makefile-ul poate folosi calea:

\begin{lstlisting}[style=code]
/usr/lib/llvm-18/include
/usr/lib/llvm-18/lib
\end{lstlisting}

\section{Compilare}

Pentru compilare:

\begin{lstlisting}[style=code]
make clean
make
\end{lstlisting}

Executabile generate:

\begin{lstlisting}[style=code]
analyzer
main
server
client
admin_client
watch_reports
\end{lstlisting}

\section{Rulare server}

Serverul se porneste cu:

\begin{lstlisting}[style=code]
./server
\end{lstlisting}

Output asteptat:

\begin{lstlisting}[style=code]
INET server running on 127.0.0.1:8080
UNIX admin socket running on /tmp/t17_admin.sock
Waiting for clients with select()...
\end{lstlisting}

Serverul asculta simultan pe:

\begin{itemize}
    \item \texttt{127.0.0.1:8080} pentru clienti normali;
    \item \texttt{/tmp/t17\_admin.sock} pentru admin.
\end{itemize}

\section{Rulare client C normal}

Pentru upload si analiza:

\begin{lstlisting}[style=code]
./client tests/sample.c
\end{lstlisting}

Pentru alt fisier:

\begin{lstlisting}[style=code]
./client tests/bad_type.c
\end{lstlisting}

Pentru download de raport:

\begin{lstlisting}[style=code]
./client download sample_report.txt
\end{lstlisting}

Rapoartele descarcate se salveaza in:

\begin{lstlisting}[style=code]
downloads/
\end{lstlisting}

\section{Rulare client Python}

Clientul Python poate fi folosit pentru upload:

\begin{lstlisting}[style=code]
python3 python_client/client.py upload tests/sample.c
\end{lstlisting}

Pentru download:

\begin{lstlisting}[style=code]
python3 python_client/client.py download sample_report.txt
\end{lstlisting}

Rapoartele descarcate prin clientul Python se salveaza in:

\begin{lstlisting}[style=code]
downloads_py/
\end{lstlisting}

\section{Rulare admin client}

Admin clientul se porneste cu:

\begin{lstlisting}[style=code]
./admin_client
\end{lstlisting}

Meniul admin:

\begin{lstlisting}[style=code]
1. Server stats
2. Show server logs
3. List uploaded files
4. List reports
5. Server status
6. Clear server logs
7. Delete report
8. List users
9. Show analysis jobs
10. Quit
\end{lstlisting}

Admin clientul comunica prin socket UNIX:

\begin{lstlisting}[style=code]
/tmp/t17_admin.sock
\end{lstlisting}

\section{Monitorizare rapoarte cu inotify}

Pentru monitorizarea directorului \texttt{reports/}:

\begin{lstlisting}[style=code]
./watch_reports
\end{lstlisting}

Apoi, intr-un alt terminal:

\begin{lstlisting}[style=code]
./client tests/sample.c
\end{lstlisting}

Output posibil:

\begin{lstlisting}[style=code]
Watching directory: reports
Waiting for report changes...
[2026-06-12 12:26:09] report created: sample_report.txt
[2026-06-12 12:26:09] report modified: sample_report.txt
[2026-06-12 12:26:09] report written: sample_report.txt
\end{lstlisting}

\section{Fisiere de configurare}

\subsection{\texttt{config/users.cfg}}

Format:

\begin{lstlisting}[style=code]
username:password:role
\end{lstlisting}

Exemplu:

\begin{lstlisting}[style=code]
user1:pass1:user
admin:admin123:admin
\end{lstlisting}

\subsection{\texttt{config/app.cfg}}

Fisier folosit pentru configurarea analyzer-ului.

\section{Directoare generate la runtime}

\begin{itemize}
    \item \texttt{uploads/} - fisiere primite de server;
    \item \texttt{reports/} - rapoarte generate;
    \item \texttt{logs/} - loguri, statistici si joburi;
    \item \texttt{downloads/} - rapoarte descarcate de clientul C;
    \item \texttt{downloads\_py/} - rapoarte descarcate de clientul Python.
\end{itemize}

\section{Testare recomandata}

\subsection{Test complet}

Terminal 1:

\begin{lstlisting}[style=code]
./server
\end{lstlisting}

Terminal 2:

\begin{lstlisting}[style=code]
./watch_reports
\end{lstlisting}

Terminal 3:

\begin{lstlisting}[style=code]
./client tests/sample.c
./client tests/bad_type.c
./client download sample_report.txt
\end{lstlisting}

Terminal 4:

\begin{lstlisting}[style=code]
./admin_client
\end{lstlisting}

In admin client se pot testa:

\begin{itemize}
    \item \texttt{Server stats};
    \item \texttt{List reports};
    \item \texttt{List users};
    \item \texttt{Show analysis jobs};
    \item \texttt{Delete report}.
\end{itemize}

\section{Curatare}

Pentru stergerea executabilelor si a directorului \texttt{uploads/}:

\begin{lstlisting}[style=code]
make clean
\end{lstlisting}

Pentru curatare manuala:

\begin{lstlisting}[style=code]
rm -rf uploads reports logs downloads downloads_py
rm -f /tmp/t17_admin.sock
\end{lstlisting}

\section{Documentatie}

Documentatia oficiala este in LaTeX:

\begin{lstlisting}[style=code]
docs/SRS_SDD.tex
docs/protocol.tex
README.tex
\end{lstlisting}

Pentru generare PDF:

\begin{lstlisting}[style=code]
pdflatex README.tex
pdflatex -output-directory=docs docs/SRS_SDD.tex
pdflatex -output-directory=docs docs/protocol.tex
\end{lstlisting}

\section{Functionalitati principale bifate}

\begin{itemize}
    \item socket INET TCP pentru clienti normali;
    \item socket UNIX pentru admin;
    \item \texttt{select()} pentru monitorizarea socket-urilor;
    \item \texttt{fork()} pentru clienti multipli;
    \item \texttt{exec()} si \texttt{pipe()} pentru rularea analyzer-ului;
    \item upload si download de fisiere;
    \item operatii admin;
    \item job tracking;
    \item \texttt{inotify} pentru monitorizare rapoarte;
    \item client in alta limba, prin Python.
\end{itemize}

\section{Concluzie}

Proiectul este o aplicatie distribuita locala care foloseste mai multe concepte importante din programarea concurenta si distribuita: socket-uri INET, socket-uri UNIX, procese, pipe-uri, select, autentificare, transfer de fisiere si monitorizare cu inotify.

\end{document}